\documentclass[a4paper,12pt]{article}
\usepackage[T2A]{fontenc}
\usepackage[utf8]{inputenc}
\usepackage[russian]{babel}
\usepackage{cmap}
\usepackage{amsmath,amssymb,mathtools}
\usepackage{geometry}
\geometry{left=25mm,right=25mm,top=24mm,bottom=24mm}
\usepackage{microtype}
\usepackage{icomma}
\usepackage{tikz}
\usetikzlibrary{arrows.meta,positioning,shapes.geometric,calc}
\usepackage{tabularx,booktabs,array}
\usepackage{enumitem}
\usepackage{xcolor}
\usepackage{fancyhdr}
\usepackage{setspace}
\usepackage{hyperref}
\definecolor{accent}{HTML}{087E82}
\definecolor{darkaccent}{HTML}{07555A}
\definecolor{textgray}{HTML}{4C4C4C}
\definecolor{faint}{HTML}{A8BFC0}
\hypersetup{hidelinks,pdftitle={Чек-лист. Основы тригонометрии},pdfauthor={Лёвин Артём Александрович}}
\setstretch{1.15}
\setlength{\parindent}{0pt}
\setlength{\parskip}{5pt}
\setlist[enumerate]{leftmargin=*,itemsep=7pt,topsep=6pt}
\pagestyle{fancy}
\fancyhf{}
\fancyhead[L]{\footnotesize\color{textgray}Основы тригонометрии}
\fancyhead[R]{\footnotesize\color{textgray}Анна Баннова}
\fancyfoot[C]{\footnotesize\color{textgray}\thepage}
\renewcommand{\headrulewidth}{0.35pt}
\renewcommand{\headrule}{\hbox to\headwidth{\color{accent}\leaders\hrule height\headrulewidth\hfill}}
\renewcommand{\footrulewidth}{0pt}
\newcommand{\sectionline}{\par\vspace{-2mm}{\color{accent}\rule{\linewidth}{0.35pt}}\par\vspace{2mm}}
\newcommand{\topic}[1]{\section*{\normalfont\color{darkaccent}#1}\sectionline}
\newcommand{\key}[1]{\textcolor{darkaccent}{#1}}
\begin{document}
\thispagestyle{empty}
\begin{center}
\vspace*{17mm}
{\Large\color{darkaccent} МАТЕМАТИКА}\par
\vspace{23mm}
{\Huge\color{textgray} Чек-лист}\par
\vspace{11mm}
{\LARGE\color{accent} Основы тригонометрии}\par
\vspace{3mm}
{\Large\color{accent} Окружность, значения и обратные функции}\par
\vspace{16mm}
{\large Подготовка к профильной математике}\par
\vfill
{\Large Анна Баннова}\par
\vspace{12mm}
{\large 22 ноября 2026 г.}\par
\vspace{18mm}
{\normalsize Лёвин Артём Александрович}\par
{\normalsize эксклюзивно для Анны Банновой}\par
\vspace{17mm}
\begin{tikzpicture}[x=1cm,y=1cm]
  \draw[accent,line width=0.8pt] (-4,0) .. controls (-2.6,0.55) and (-1.25,-0.47) .. (0,0)
     .. controls (1.25,0.50) and (2.55,-0.52) .. (4,0);
  \fill[accent] (0,0) circle (1.6pt);
\end{tikzpicture}
\end{center}
\clearpage

\topic{1. Единичная окружность}
Единичная окружность --- окружность радиуса $1$ с центром в начале координат. Угол $\alpha$ откладывают от положительного направления оси $x$. \key{Положительный угол} отсчитывают против часовой стрелки, \key{отрицательный} --- по часовой.

\textbf{Главная связь:} точка, соответствующая углу $\alpha$, имеет координаты
\[
\boxed{P_\alpha=(\cos\alpha\,;\,\sin\alpha).}
\]
Поэтому косинус показывает \emph{горизонтальную} координату, а синус --- \emph{вертикальную}.

\begin{center}
\begin{tikzpicture}[scale=1.25,>=Latex]
 \draw[->,thin,gray!80] (-3.2,0)--(3.45,0) node[right,black] {$x$};
 \draw[->,thin,gray!80] (0,-2.75)--(0,2.95) node[above,black] {$y$};
 \draw[accent,line width=1.1pt] (0,0) circle (2.35);
 \coordinate (P) at (1.662,1.662);
 \draw[accent,thick,->] (0,0)--(P);
 \draw[accent,dashed] (P)--(1.662,0);
 \draw[accent,dashed] (P)--(0,1.662);
 \fill[accent] (P) circle (2.1pt);
 \draw[darkaccent,->] (0.75,0) arc[start angle=0,end angle=45,radius=0.75];
 \node[darkaccent] at (1.05,0.43) {$\alpha$};
 \node[anchor=south west] at (1.68,1.72) {$P_\alpha$};
 \node[below] at (1.66,-0.05) {$\cos\alpha$};
 \node[left] at (-0.05,1.65) {$\sin\alpha$};
 \node[below right] at (2.35,0) {$A(1;0)$};
 \node[above right] at (0,2.35) {$B(0;1)$};
 \node[below left] at (-2.35,0) {$C(-1;0)$};
 \node[below left] at (0,-2.35) {$D(0;-1)$};
 \node[below left] at (0,0) {$O$};
\end{tikzpicture}
\end{center}
\textbf{Знаки по четвертям:} в I четверти оба значения положительны; во II $\sin\alpha>0$, $\cos\alpha<0$; в III оба отрицательны; в IV $\sin\alpha<0$, $\cos\alpha>0$.

\topic{2. Градусы и радианы}
\[
180^\circ=\pi\ \text{рад},\qquad
\alpha_{\text{рад}}=\alpha_{^\circ}\cdot\frac{\pi}{180^\circ},\qquad
\alpha_{^\circ}=\alpha_{\text{рад}}\cdot\frac{180^\circ}{\pi}.
\]
\textbf{Примеры:}
\[
45^\circ=\frac{45\pi}{180}=\frac\pi4,\qquad
120^\circ=\frac{120\pi}{180}=\frac{2\pi}{3},\qquad
-90^\circ=-\frac\pi2.
\]
\textbf{Опорные углы:} $30^\circ=\frac\pi6$, $45^\circ=\frac\pi4$, $60^\circ=\frac\pi3$, $90^\circ=\frac\pi2$.

\topic{3. Табличные значения}
\begingroup\small
\renewcommand{\arraystretch}{1.55}
\begin{center}
\begin{tabular}{>{\raggedright\arraybackslash}p{20mm}*{5}{>{\centering\arraybackslash}p{22mm}}}
\toprule
Угол & $0$ & $\frac\pi6$ & $\frac\pi4$ & $\frac\pi3$ & $\frac\pi2$ \\
\midrule
$\sin\alpha$ & $0$ & $\frac12$ & $\frac{\sqrt2}{2}$ & $\frac{\sqrt3}{2}$ & $1$ \\
$\cos\alpha$ & $1$ & $\frac{\sqrt3}{2}$ & $\frac{\sqrt2}{2}$ & $\frac12$ & $0$ \\
$\tg\alpha$ & $0$ & $\frac{\sqrt3}{3}$ & $1$ & $\sqrt3$ & --- \\
$\ctg\alpha$ & --- & $\sqrt3$ & $1$ & $\frac{\sqrt3}{3}$ & $0$ \\
\bottomrule
\end{tabular}
\end{center}
\endgroup
Знак «---» означает, что значение \textbf{не определено}.
\[
\boxed{\tg\alpha=\frac{\sin\alpha}{\cos\alpha}\quad(\cos\alpha\ne0),\qquad
\ctg\alpha=\frac{\cos\alpha}{\sin\alpha}\quad(\sin\alpha\ne0).}
\]
\textbf{Пример:} $\displaystyle\tg\frac\pi6=\frac{1/2}{\sqrt3/2}=\frac1{\sqrt3}=\frac{\sqrt3}{3}.$

\topic{4. Обратные тригонометрические функции}
\textbf{Задача обратной функции} --- найти угол по известному значению синуса, косинуса, тангенса или котангенса. Чтобы ответ был единственным, у каждой функции установлен свой \key{основной промежуток}.

\begingroup\small
\renewcommand{\arraystretch}{1.5}
\begin{center}
\begin{tabularx}{\linewidth}{@{}l X X@{}}
\toprule
Функция & Что означает $\alpha$ & Основной промежуток \\
\midrule
$\alpha=\arcsin t$ & $\sin\alpha=t$, $-1\le t\le1$ & $\left[-\frac\pi2;\frac\pi2\right]$ \\
$\alpha=\arccos t$ & $\cos\alpha=t$, $-1\le t\le1$ & $[0;\pi]$ \\
$\alpha=\arctg t$ & $\tg\alpha=t$, $t\in\mathbb R$ & $\left(-\frac\pi2;\frac\pi2\right)$ \\
$\alpha=\arcctg t$ & $\ctg\alpha=t$, $t\in\mathbb R$ & $(0;\pi)$ \\
\bottomrule
\end{tabularx}
\end{center}
\endgroup
\textbf{Важно:} для $\arcctg$ используется школьное соглашение с промежутком $(0;\pi)$.

\textbf{Примеры поиска по таблице и основному промежутку:}
\[
\arcsin\frac12=\frac\pi6,\qquad
\arccos\frac{\sqrt3}{2}=\frac\pi6,\qquad
\arctg1=\frac\pi4.
\]
Если нужен угол из промежутка $(-\pi/2;\pi/2)$, для которого $\tg\alpha=10$, точный ответ записывают как $\alpha=\arctg10$.

\topic{5. Отрицательный аргумент: четыре правила}
Для допустимых положительных $t$ действуют полезные соотношения:
\begin{align*}
\arcsin(-t)&=-\arcsin t,\\
\arccos(-t)&=\pi-\arccos t,\\
\arctg(-t)&=-\arctg t,\\
\arcctg(-t)&=\pi-\arcctg t.
\end{align*}
Особенно внимательно различайте \emph{минус перед ответом} и \emph{вычитание ответа из $\pi$}.

\textbf{Четыре примера из одного приёма:}
\begin{align*}
\arcsin\!\left(-\frac{\sqrt3}{2}\right)&=-\frac\pi3,\\[2pt]
\arccos\!\left(-\frac{\sqrt2}{2}\right)&=\pi-\frac\pi4=\frac{3\pi}{4},\\[2pt]
\arctg(-\sqrt3)&=-\frac\pi3,\\[2pt]
\arcctg(-\sqrt3)&=\pi-\frac\pi6=\frac{5\pi}{6}.
\end{align*}
\textbf{Контроль:} $\arccos(-1/2)=2\pi/3$, поскольку $2\pi/3$ находится во II четверти и его косинус равен $-1/2$.

\topic{6. Если функция и обратная функция стоят рядом}
Нельзя механически «сокращать» $\arccos(\cos\beta)$ и $\arcsin(\sin\beta)$ для произвольного $\beta$. Сначала вычислите значение прямой функции, затем найдите угол из \textbf{основного промежутка}.

\textbf{Пример 1.}
\[
\arccos\left(\cos\left(-\frac\pi3\right)\right)
=\arccos\frac12=\frac\pi3.
\]
\textbf{Пример 2.}
\[
\arcsin\left(\sin\frac{2\pi}{3}\right)
=\arcsin\frac{\sqrt3}{2}=\frac\pi3.
\]
Второй результат равен $\pi/3$, поскольку угол $2\pi/3$ лежит за пределами основного промежутка арксинуса.

\clearpage
\thispagestyle{fancy}
\begin{center}
{\Large\color{darkaccent} Алгоритм: как найти обратный тригонометрический угол}\par
\vspace{13mm}
\begin{tikzpicture}[node distance=12mm,>=Latex, every node/.style={align=center,font=\small},
 st/.style={ellipse,draw=accent,line width=0.75pt,text width=40mm,minimum height=11mm},
 ac/.style={rectangle,rounded corners=2pt,draw=darkaccent,line width=0.65pt,text width=99mm,minimum height=13mm,inner sep=3mm},
 dec/.style={diamond,aspect=2.5,draw=darkaccent,line width=0.7pt,text width=27mm,inner sep=1.2mm},
 branch/.style={rectangle,rounded corners=2pt,draw=darkaccent,line width=0.65pt,text width=53mm,minimum height=22mm,inner sep=2mm},
 arrow/.style={-Latex,semithick,draw=darkaccent}]
 \node[st] (start) {Начало};
 \node[ac,below=of start] (range) {Установите обратную функцию и основной промежуток; если аргумент недопустим, выражение не определено};
 \node[dec,below=15mm of range] (sign) {Аргумент отрицателен?};
 \node[branch,below left=17mm and 8mm of sign] (minus) {Найдите угол для модуля аргумента и примените правило: \( -\alpha \) или \(\pi-\alpha\)};
 \node[branch,below right=17mm and 8mm of sign] (plus) {Найдите табличный угол для положительного аргумента};
 \node[ac,below=65mm of sign] (choose) {Выберите угол, принадлежащий основному промежутку данной обратной функции};
 \node[ac,below=of choose] (check) {Проверьте ответ подстановкой в синус, косинус, тангенс или котангенс};
 \node[st,below=of check] (end) {Ответ};
 \draw[arrow] (start)--(range);
 \draw[arrow] (range)--(sign);
 \draw[arrow] (sign.south west)--node[above left,font=\scriptsize]{да}(minus.north);
 \draw[arrow] (sign.south east)--node[above right,font=\scriptsize]{нет}(plus.north);
 \coordinate (merge) at ([yshift=10mm]choose.north);
 \draw[darkaccent,semithick] (minus.south) |- (merge);
 \draw[darkaccent,semithick] (plus.south) |- (merge);
 \draw[arrow] (merge) -- (choose.north);
 \draw[arrow] (choose)--(check);
 \draw[arrow] (check)--(end);
\end{tikzpicture}
\end{center}
\clearpage

\topic{7. Типичные ошибки и способы проверки}
\begin{enumerate}[label=\textcolor{accent}{\arabic*.},itemsep=4pt]
\item \textbf{Путают координаты точки:} на окружности $x=\cos\alpha$, $y=\sin\alpha$.
\item \textbf{Забывают направление отсчёта:} отрицательный угол откладывается по часовой стрелке.
\item \textbf{Считают $\arccos(-t)=-\arccos t$:} ответ арккосинуса обязан принадлежать $[0;\pi]$.
\item \textbf{Считают $\arcctg(-t)=-\arcctg t$:} в принятом соглашении результат принадлежит $(0;\pi)$.
\item \textbf{Пишут $\tg(\pi/2)=\infty$:} правильная запись --- \emph{значение не определено}, так как знаменатель равен нулю.
\item \textbf{Сокращают обратные функции без проверки интервала:} $\arcsin(\sin(2\pi/3))=\pi/3$.
\end{enumerate}
\textbf{Короткая памятка:} определите функцию $\rightarrow$ проверьте область значений аргумента $\rightarrow$ вспомните основной промежуток $\rightarrow$ найдите угол $\rightarrow$ выполните проверку.

\clearpage
\topic{Тренировочный блок}
\textit{Путь А. Десять упражнений от базового уровня к повышенному. Ответы помещены на следующей странице.}
\begin{enumerate}[label=\textcolor{accent}{\bfseries\arabic*.},itemsep=8pt]
 \item Переведите в радианы: $135^\circ$, $225^\circ$, $-60^\circ$.
 \item Переведите в градусы: $\dfrac{7\pi}{6}$, $-\dfrac{3\pi}{4}$, $\dfrac{5\pi}{3}$.
 \item Точка единичной окружности имеет координаты $P\left(-\dfrac{\sqrt3}{2};\dfrac12\right)$. Найдите $\cos\alpha$, $\sin\alpha$ и укажите четверть.
 \item Вычислите: $\sin\dfrac{2\pi}{3}$ и $\cos\dfrac{2\pi}{3}$.
 \item Вычислите: $\sin\left(-\dfrac\pi6\right)$ и $\cos\left(-\dfrac\pi6\right)$.
 \item Вычислите: $\tg\dfrac\pi6$ и $\ctg\dfrac\pi3$.
 \item Для каждого выражения укажите значение либо запишите «не определено»: $\tg\dfrac\pi2$, $\ctg0$, $\tg0$, $\ctg\dfrac\pi2$.
 \item Найдите: $\arcsin\left(-\dfrac{\sqrt2}{2}\right)$ и $\arctg1$.
 \item Найдите: $\arccos\left(-\dfrac12\right)$ и $\arcctg(-\sqrt3)$.
 \item Вычислите:
 \[
 \arccos\left(-\frac{\sqrt2}{2}\right)
 +\arcsin\left(-\frac{\sqrt3}{2}\right)
 +\arctg(-\sqrt3)+\arcctg(-\sqrt3).
 \]
\end{enumerate}
\clearpage
\topic{Ответы}
\begingroup\renewcommand{\arraystretch}{1.5}
\begin{center}
\begin{tabularx}{\linewidth}{@{}c X@{}}
\toprule
№ & Ответ \\
\midrule
1 & $\dfrac{3\pi}{4}$; $\dfrac{5\pi}{4}$; $-\dfrac\pi3$. \\
2 & $210^\circ$; $-135^\circ$; $300^\circ$. \\
3 & $\cos\alpha=-\dfrac{\sqrt3}{2}$, $\sin\alpha=\dfrac12$; II четверть. \\
4 & $\dfrac{\sqrt3}{2}$; $-\dfrac12$. \\
5 & $-\dfrac12$; $\dfrac{\sqrt3}{2}$. \\
6 & $\dfrac{\sqrt3}{3}$; $\dfrac{\sqrt3}{3}$. \\
7 & Не определено; не определено; $0$; $0$. \\
8 & $-\dfrac\pi4$; $\dfrac\pi4$. \\
9 & $\dfrac{2\pi}{3}$; $\dfrac{5\pi}{6}$. \\
10 & $\dfrac{11\pi}{12}$. \\
\bottomrule
\end{tabularx}
\end{center}
\endgroup
\vspace{9mm}
\textbf{Итоговая самопроверка.} Если в задачах 8--10 получился угол вне основного промежутка соответствующей обратной функции, вернитесь к таблице из раздела 4. Проверьте знаки и сложение дробей с $\pi$.

\vfill
\begin{center}
{\small\color{textgray} Лёвин Артём Александрович --- эксклюзивно для Анны Банновой}
\end{center}
\end{document}
